\subsubsection{Linear cohomology bounds for the original tame boundary}

We retain the original boundary throughout this subsection.  Let $S$ be
a normal proper connected surface over an algebraically closed field $k$
of characteristic $p>0$, let $T\subsetneq S$ be closed, and put
$U=S-T$.  Write $D$ for the reduced divisorial part of $T$.  We denote by
$Q(S,T)$ the quotient of $\pi_1(U)$ obtained by killing the wild inertia
at the original generic points of $D$.  Thus no condition is imposed
at the additional exceptional valuations of a resolution.

For a finite-dimensional $\mathbf F_p[Q(S,T)]$-module $M$, of dimension
$r$, choose a finite Galois cover of $U$ which trivializes $M$, normalize
$S$ in that cover, and write $f:S'\to S$ and $G$ for the resulting
finite map and group.  We use the canonical lattice
\[
             E=(f_*\mathcal O_{S'}\otimes_{\mathbf F_p}M)^G.
                                                        \tag{D.TC.1}
\]
This definition is independent of the chosen trivializing cover.
It defines a coherent reflexive sheaf of rank $r$.  Indeed, the
normal finite surface above $S$ is Cohen--Macaulay; kernels between
reflexive sheaves with torsion-free target are reflexive.  In (D.TC.1)
we use invariants only as a kernel.  We do not assume that the order
of $G$ is invertible, or that taking $G$-invariants is exact.

The ordinary Frobenius map induces a generically invertible injection
\[
            \varphi:F_S^{[*]}E\longrightarrow E.          \tag{D.TC.2}
\]
At a tame boundary DVR it is the invariant map on a Kummer lattice.
If a character has weight $j/n$, where $0\le j<n$ and $(n,p)=1$,
its elementary exponent is
\[
                         c_j=\lfloor pj/n\rfloor,
                   \qquad 0\le c_j\le p-1.              \tag{D.TC.3}
\]
Away from $D$, (D.TC.2) is an isomorphism.  Its cokernel
$\mathcal T$ is pure of dimension one: both terms in (D.TC.2)
have depth two at closed points, and the depth lemma excludes
zero-dimensional associated points in the cokernel.

\begin{lemma}[Proved: the summed boundary-lattice index]
\label{D:tame-boundary-summed-index}
Fix a completed normal surface germ $R$ at one of the finitely many
singular or boundary-intersection points.  Choose a reduced element
$a\in R$ whose divisor contains the local original boundary, and a
parameter $x$ such that
\[
                    A_0=k[[x,a]]\subset R
\]
is finite and $x$ is separating on every component of the reduced
curve $R/(a)$.  There is a constant $L$, depending only on these
fixed data, with the following property.

For every tame local system of rank $r$, and every prime-to-$p$
Kummer denominator $n$ which kills its original divisorial inertia,
the pure cokernel $\mathcal T$ of (D.TC.2) admits, as a
$k[[x]]$-lattice, a comparison
\[
         T_{\rm diag}\ \subset\ \mathcal T\
                      \subset\ x^{-L}T_{\rm diag},
       \qquad
         \operatorname{rank}_{k[[x]]}T_{\rm diag}
                      \le (p-1)[R:A_0]r.               \tag{D.TC.4}
\]
The generic graded pieces of $T_{\rm diag}$ are the whole
character eigenspaces, repeated $c_j$ times as in (D.TC.3).
For any one generic filtration of these pieces, the sum of the
positive index corrections after passage to the normalization
of the fixed curve $R/(a)$ is at most $C r$, with $C$ independent
of $n$, of $M$, and of the number of characters.
\end{lemma}

\begin{proof}
Put $A_n=k[[x,s]]$, with $a=s^n$, and
$R_n=R\otimes_{A_0}A_n$.  The reduced-divisor condition and
$(n,p)=1$ make $R_n$ normal in codimension one; it is
Cohen--Macaulay because it is finite free over $R$.  Hence it is
normal.  At the closed point it is local, with residue field $k$.
The pulled-back local system extends over the puncture of $R_n$,
by the regular codimension-one Kummer calculation and normality.
Let $E_n$ be its reflexive unit lattice.  It is free over the
regular ring $A_n$, of rank $d_0r$, where $d_0=[R:A_0]$.

We use the fixed stable-MCM factorization ideal from
Lemma~\ref{D:Kummer-uniform-Frobenius-defect} and
Theorem~\ref{D:local-tame-p-linear-H1}.  Their local
tame-cohomology argument supplies a fixed
power $x^N$ whose Frobenius power annihilates the ordinary
Frobenius defect of every $E_n$, uniformly in $n$ and $r$.
This is the uniform ideal obtained from the base Jacobian
minors and the term $a\operatorname{Jac}(R)$; it is not an
assertion that test ideals commute with Kummer base change.
Indeed, the fixed Jacobian ideal is contained in the sum of
Noether differents.  Each such element factors the identity
on an MCM module through a free module, and its $p$th power
annihilates the ordinary-to-reflexive Frobenius defect
\citep[section~2; section~3, formula~(3.1) and
Examples~3.2--3.3, pp.~2--3 and 6--7]{C_IyengarTakahashiStable}.
Equivalently, after increasing a fixed integer $K$, the cokernel
of $F_{R_n}^*E_n\to E_n$ is killed by $x^K$.

Consider relative Frobenius for the finite free $A_0$-algebra $R$:
\[
             \alpha:F_{A_0}^*R\longrightarrow R.
\]
Its determinant $h(x,a)$ is not zero on the generic point of
any component of $a=0$.  This is exactly where the separating
choice of $x$ is used: the corresponding finite extension of
curve function fields is separable, so relative Frobenius is
an isomorphism there.  Thus $h(x,0)\ne0$.  The adjugate
annihilates the relative cokernel.  The same assertion after
$a=s^n$ gives $h(x,s^n)$ for $R_n/A_n$.

Let $B$ be the matrix, between free $A_n$-modules of rank
$d_0r$, of the composite
\[
       F_{A_n}^*E_n\longrightarrow F_{R_n}^*E_n
                              \longrightarrow E_n.
\]
Tensoring the relative-Frobenius cokernel calculation and
then using the ordinary defect yields
\[
                 x^K h(x,s^n)\operatorname{coker}(B)=0.
                                                        \tag{D.TC.5}
\]
Consequently $x^K h B^{-1}$ has integral entries.
Put $h_0=\operatorname{ord}_x h(x,0)$.  Modulo $s^{np}$,
the inverse of $h(x,s^n)$ is the finite geometric expansion
of its constant term in $s^n$, with at most $p$ terms.
Its $x$-denominator is bounded by $x^{p h_0}$, up to a unit.
Thus one fixed $L$, for example $K+p h_0$ after enlarging
it to absorb the chosen unit convention, satisfies
\[
 B^{-1}\bmod s^{np}\ \in\
     x^{-L}\operatorname{Mat}_{d_0r}
                        (A_n/(s^{np})).                 \tag{D.TC.6}
\]
The bound is independent of the size of the matrix.

Choose a homogeneous $A_n$-basis for $E_n$ under $\mu_n$.
It exists by choosing a homogeneous basis modulo $(x,s)$
and applying Nakayama.  If a basis vector has weight $-j$,
the corresponding invariant generator is $s^jv_j$.
The invariant module
$(F_{A_n}^*E_n)^{\mu_n}$ has generator
$s^{pj\bmod n}Fv_j$, and the map from $F_{A_0}^*E$
to it is diagonal with entries $a^{c_j}$.
The quotient of this diagonal map is the lattice
$T_{\rm diag}$ in the statement.  It is free over $k[[x]]$
and has rank $\sum c_j\le(p-1)d_0r$.

The invariant map induced by $B$ sends the diagonal image
into the ordinary Frobenius image in $E$, hence into the
reflexive image in (D.TC.2).  It therefore gives
$T_{\rm diag}\to\mathcal T$.  This map is an isomorphism
after inverting $x$.  Purity makes $\mathcal T$ torsion-free
over $k[[x]]$, while $T_{\rm diag}$ is free, so it is
injective.  Also $a^{p-1}\mathcal T=0$, since this holds at
each generic boundary point and $\mathcal T$ has no
zero-dimensional submodule.  Equation (D.TC.6), applied
modulo $a^p$, now gives
$x^L\mathcal T\subset T_{\rm diag}$.  This proves (D.TC.4).

For completeness, the final-grade estimate uses a single
filtration, not a separate index bound for each character.
Identify the two lattices over $k((x))$, choose the generic
filtration by boundary component, character and $a$-layer,
and intersect it with both lattices.  Additivity of length
in the resulting finite filtrations gives
\[
 \sum_{\rm grades}\operatorname{length}_{k[[x]]}
        (L_{\rm grade}/L'_{\rm grade})
     =\operatorname{length}_{k[[x]]}(\mathcal T/T_{\rm diag})
     \le L(p-1)d_0r.                                    \tag{D.TC.7}
\]
Let $\overline R_0$ be the normalization of $R/(a)$.
It is finite free over $k[[x]]$, of fixed rank.  For an
inclusion $L'\subset L$ of the lattices occurring in
(D.TC.7), the quotient of their normalized spans is a
quotient of
\[
             \overline R_0\otimes_{k[[x]]}(L/L').
\]
Thus normalization multiplies the total index by at most
$\operatorname{rank}_{k[[x]]}\overline R_0$.  There is no
additional factor equal to the number of grades or to $r$.
This proves the last assertion.
\end{proof}

\begin{lemma}[Proved: a uniform Euler bound for the tame Frobenius quotient]
\label{D:tame-canonical-Euler}
There is a constant $C$, depending only on $(S,T)$, such that
the pure quotient in (D.TC.2) satisfies
\[
                              \chi(S,\mathcal T)\le Cr.  \tag{D.TC.8}
\]
The same bound holds at every stage of a finite cyclic
coefficient-Frobenius orbit of canonical lattices.
\end{lemma}

\begin{proof}
Enlarge a fixed finite set $\Sigma$ so that $S-\Sigma$ is
smooth and its boundary components are smooth and disjoint.
On this open set the assertion is the elementary Kummer
calculation.  A character of weight $j/n$ contributes the
whole weight space in exactly $c_j$ successive normal
layers.  The possible normal twists belong to the fixed
list $0,\ldots,p-2$.

We explain both the global slope control and the correction
at $\Sigma$.  Normalize each fixed boundary curve, obtaining
smooth proper curves $\overline D_i$.  The normal line
bundle on $D_i-\Sigma$ has a fixed extension to
$\overline D_i$.  Choose finitely many fixed rational
normal trivializations.  Kummer root changes between
them differ by roots of these fixed rational functions.
For a weight $j/n\in[0,1)$, the corresponding integral
lattice shift is bounded by the absolute orders of
these functions, independently of $j$ and $n$.

At a point of $\Sigma$, use the local cover $R_n$ of the
preceding lemma.  The module $E_n/sE_n$ is a
maximal-Cohen--Macaulay module on the reduced curve
$R/(a)$.  If $b$ is a non-zero-divisor in its fixed
conductor, then
\[
 b\bigl(\overline R_0\otimes(E_n/sE_n)/{\rm torsion}\bigr)
                  \subset E_n/sE_n
                  \subset
       \overline R_0\otimes(E_n/sE_n)/{\rm torsion}.
                                                        \tag{D.TC.9}
\]
Taking character summands is exact here, since
$(n,p)=1$.  The same fixed conductor and the preceding
normal-trivialization shifts therefore glue the normalized
character lattices to vector bundles $P_{i,j}$ on
$\overline D_i$, with comparison divisors independent
of $n,r$ and $j$.  Equivalently, one may carry out the
gluing on the root gerbe of the chosen normal line bundle
and then twist the character-$j$ part by the universal
character line before descent
\citep[Proposition~3 and Theorem~4, pp.~11--16]{Borne2007D}
\citep[Theorem~6.1, pp.~33--36]{BorneVistoli2012D}.
Formula (D.TC.9) controls
the difference from the chosen normalized local lattice.

The regular Frobenius map on $E_n$ gives, after reducing
modulo $s$, normalizing and accounting for these fixed
comparison divisors, generically invertible maps
\[
       F_{\overline D_i}^*P_{i,j}
                  \longrightarrow
       P_{i,pj\bmod n}(B_i+c_jN_i).                      \tag{D.TC.10}
\]
Here $B_i$ is fixed and $N_i$ denotes the chosen normal
line bundle; changing its sign convention merely changes
the fixed bound.  Regularity of this map uses only the
forward Frobenius map.  We do not assert an inverse
without a pole bound.  The finitely many comparison
divisors in (D.TC.9), and the bound $0\le c_j\le p-1$,
give one constant $C_i$ such that
\[
 p\,\mu_{\max}(P_{i,j})
             \le \mu_{\max}(P_{i,pj\bmod n})+C_i.
\]
Cycling the characters gives
$\mu_{\max}(P_{i,j})\le C_i/(p-1)$.
For coefficient fields $\mathbf F_{p^f}$ the same argument
cycles pairs consisting of a coefficient-Frobenius stage
and a character.  It is unchanged and remains independent
of $f$.

Thus the Euler characteristic of each whole character
bundle, with any of the finitely many normal twists above,
is at most a fixed constant times its rank.  The sum of
these ranks over all normal layers is at most
$(p-1)r$ for each fixed $D_i$.

The actual pure grades of $\mathcal T$ need not equal
these normalized character lattices at the points of
$\Sigma$.  This is precisely the correction estimated
in Lemma~\ref{D:tame-boundary-summed-index}.  On the
generic curve the diagonal source grades are whole
character spaces, since $c_j$ is constant on a character;
the leading term of $B$ maps them integrally into the
target.  Intersect the same generic filtration with
$T_{\rm diag}$ and $\mathcal T$.  Equation (D.TC.7) and
the normalized-span estimate give a total positive
Euler correction $O(r)$, not one $O(r)$ correction for
each grade.  Summing over the fixed finite set $\Sigma$
proves (D.TC.8).
\end{proof}

\begin{proposition}[Proved: linear coherent cohomology for canonical tame lattices]
\label{D:tame-canonical-linear-cohomology}
For each $i=0,1,2$ there is a constant $C_i$, depending
only on $(S,T)$, such that every canonical lattice (D.TC.1)
of rank $r$ satisfies
\[
                              h^i(S,E)\le C_i r.         \tag{D.TC.11}
\]
\end{proposition}

\begin{proof}
The determinant of (D.TC.2) gives
\[
       (1-p)c_1(E)=[\mathcal T]_1
                    =\sum_i m_iD_i,\qquad
                    0\le m_i\le(p-1)r.                 \tag{D.TC.12}
\]
Consequently $K_S\cdot c_1(E)=O(r)$ in the numerical
intersection theory of the normal proper surface.
By the normal-surface Riemann--Roch estimate and Frobenius
compatibility of the Chern character
\citep[Theorem~0.1]{C_Langer2026}
\citep[Theorem~4.5(2)]{C_Langer2025},
\[
 \begin{split}
 \chi(\mathcal T)
   &=\chi(E)-\chi(F_S^{[*]}E)\\
   &=-(p^2-1)\operatorname{ch}_2(E)
                  +\frac{p-1}{2}K_S\cdot c_1(E)+O(r).
 \end{split}                                            \tag{D.TC.13}
\]
The local Riemann--Roch error is bounded by a fixed
constant times the rank.  Combining (D.TC.8),
(D.TC.12), and (D.TC.13) gives
$\operatorname{ch}_2(E)\ge-Cr$, hence
\[
                                  \chi(E)\ge-Cr.         \tag{D.TC.14}
\]
The direction of this inequality is important: an upper
bound for $\chi(\mathcal T)$ supplies the lower bound
for $\chi(E)$ needed to control $h^1(E)$.

We now bound $h^0$ and $h^2$ without assuming that $S$
is projective.  Choose one smooth projective resolution
$\pi:Y\to S$ principalizing the following fixed ideal.
Codimension-one checking in (D.TC.3), followed by
reflexivity, gives
\[
    \mathcal O_S(-(p-1)D)\,E
                       \subset \varphi(F_S^{[*]}E).
\]
Multiply by the fixed stable-MCM Frobenius-defect ideal
\citep[Theorem~1.1 and Remark~2.5(1)]{C_Kimura2025}
\citep[Lemma~1.8]{C_Langer2026}
to land in the image of the ordinary $F_S^*E$.
This works for every reflexive $E$ of every rank.  The
product is one coherent ideal $I$, independent of $E$.
Write $I\mathcal O_Y=\mathcal O_Y(-A)$ and
$F=\pi^{[*]}E$.  Frobenius is flat on the smooth surface
$Y$, and the ordinary map and its fixed-ideal inclusion
give
\[
                     F(-A)\subset F_Y^*F\subset F.       \tag{D.TC.15}
\]
The second map is generically an isomorphism, hence
injective.  The first is checked at codimension-one
DVRs and then extended by reflexivity.

Choose a fixed smooth very ample curve $C$ on $Y$ which
is not contained in the determinant-zero divisors of
these maps.  Intersections with $A$ are allowed.
Both restricted maps are generically injective between
vector bundles on $C$, hence injective.  The usual
pullback inequality for maximal slopes, applied to
(D.TC.15) and its dual, yields
\[
 \mu_{\max,H}(F)\le0,\qquad
 \mu_{\max,H}(F^\vee)\le A\cdot H/(p-1),
                                                        \tag{D.TC.16}
\]
and the identical bounds on $C$, with $H$ replaced by
$C$.  Here $H=C$ may be fixed once and for all.
The maximal-slope subsheaf pulls back to a subsheaf,
so the only Frobenius slope inequality used in
(D.TC.16) is $p\mu_{\max}\le\mu_{\max}(F^*(-))$.

The vanishing $H^0(Y,F(-C))=0$ and evaluation at one
point on $C$ give
\[
                             h^0(S,E)\le h^0(Y,F)\le r.
\]
For $h^2$, let $W=\pi^{[*]}\omega_S$, a fixed line bundle.
Proper duality and reflexive pullback give an injection
\[
 H^2(S,E)^\vee=\operatorname{Hom}_S(E,\omega_S)
                       \hookrightarrow H^0(Y,F^\vee W).
                                                        \tag{D.TC.17}
\]
Put $\beta=A\cdot C/(p-1)+W\cdot C$.
Choose fixed positive integers $m,t$ with
$mC^2>\beta$ and $t>\beta$.
Then $H^0(Y,F^\vee W(-mC))=0$.  On each of the
$m$ restrictions to $C$, evaluation at $t$ distinct
points bounds the space of sections by $tr$,
using the maximal-slope bound $\beta$.
The restriction sequences therefore give
$h^2(S,E)\le mtr$.  Finally
$h^1=h^0+h^2-\chi$ and (D.TC.14) prove (D.TC.11).
\end{proof}

\begin{theorem}[Proved: all-characteristic linear degree-one cohomology for a surface boundary]
\label{D:tame-surface-all-characteristics-H1}
\label{D:surface-original-boundary-p-H1}
\label{D:surface-original-boundary-prosolvable}
For the original fixed-divisor tame quotient $Q=Q(S,T)$
there is a constant $C_{S,T}$ such that, for every finite
field $\Lambda$ and every finite-dimensional continuous
$\Lambda[Q]$-module $M$,
\[
                  \dim_\Lambda H^1(Q,M)
                              \le C_{S,T}\dim_\Lambda M.
                                                        \tag{D.TC.18}
\]
For an open subgroup $H\subset Q$ the bound is
\[
        \dim_\Lambda H^1(H,M)
                    \le C_{S,T}[Q:H]\dim_\Lambda M.       \tag{D.TC.19}
\]
In particular the maximal prosolvable quotient of $Q$
is topologically finitely generated.  We make no
finite-generation assertion here for the full group $Q$.
\end{theorem}

\begin{proof}
We first use $\Lambda=\mathbf F_p$.  On the small
\'etale site of $S$, the canonical semilinear Frobenius
gives an exact sequence
\[
          0\longrightarrow j_*M\longrightarrow E
                   \xrightarrow{\varphi-1}E
                   \longrightarrow0,                  \tag{D.TC.20}
\]
where $j:U\hookrightarrow S$.
For surjectivity, choose local generators of $E$ and
write Frobenius as a matrix on these generators.
The corresponding Lang polynomial has the form
$Cz^{[p]}-z-b$ and Jacobian $-I$.  Over the separably
closed residue field it is surjective: it is an
\'etale additive map of affine space, whose image is
an open subgroup, hence the whole connected group.
Hensel lifting at a strict stalk supplies a solution.
No freeness of $E$ at the singular point is required.
For the kernel, use (D.TC.1).  Frobenius-fixed functions
on the finite normal cover are locally constant
$\mathbf F_p$-valued functions; these extend across
the boundary.  Taking kernels and then $G$-invariants
identifies the result with $j_*M$, without an
exactness assertion about invariants.

The vector space $H^0(S,E)$ is finite dimensional and
$\varphi-1$ is surjective on it by the same Lang argument.
Hence (D.TC.20) injects $H^1(S,j_*M)$ into
$H^1(S,E)^{\varphi=1}$.  Frobenius-fixed vectors which
are independent over $\mathbf F_p$ are independent
over $k$: a shortest normalized $k$-linear relation,
minus its Frobenius transform, forces all coefficients
to lie in $\mathbf F_p$.  Proposition
\ref{D:tame-canonical-linear-cohomology} therefore gives
\[
                      \dim_{\mathbf F_p}H^1(S,j_*M)\le Cr.
                                                        \tag{D.TC.21}
\]

Apply the degree-one Leray sequence for $j$ only to
classes inflated from $H^1(Q,M)$.  These are represented
by finite torsors whose combined monodromy still factors
through the original tame quotient.  At a smooth point
of a smooth boundary divisor their local classes vanish:
the strict-local tame inertia has order prime to $p$,
and its positive cohomology on an $\mathbf F_p$-module
vanishes.  Thus their local images are supported on the
fixed finite set $\Sigma$ of singular, intersection and
isolated omitted points.

At each point of $\Sigma$, the uniform local tame-$p$
image bound in Theorem~\ref{D:local-tame-p-linear-H1}
applies to exactly these classes.  Its
Kummer cover kills the prime-to-$p$ inertia of the
given representation and torsor; the fixed Jacobian
test ideal and equivariant Matlis bound are independent
of its denominator.  The sum of the local dimensions
is therefore $O(r)$.  At a point away from the
divisorial boundary this is the already proved
punctured-surface local bound
(Theorem~\ref{D:local-p-linear-H1}).  We do not apply this
argument to all of $R^1j_*M$, which can also contain
wild classes excluded by $Q$.

Inflation injects $H^1(Q,M)$ into $H^1(U,M)$.
Combining the restricted Leray sequence, (D.TC.21),
and these finitely many local bounds proves (D.TC.18)
in characteristic $p$.  Restriction of scalars gives
the same constant for every finite field of
characteristic $p$.

For $\operatorname{char}\Lambda\ne p$, the uniform
original-boundary conductor theorem, its curvewise
consequence, and the uniform Betti estimate already
proved in Proposition~\ref{D:surface-curvewise-conductor}
give (D.TC.18).  Shapiro applied to
$\operatorname{Ind}_H^Q M$ proves (D.TC.19).
Finally the previously proved cohomology/crown
criterion, restricted to finite solvable quotients,
uses only abelian crowns and gives a uniform bound
on their number of generators.  Compactness then
gives finite generation of the maximal prosolvable
quotient.
\end{proof}

\begin{remark}[The remaining nonabelian gate]
The estimates above control all abelian crowns, in
every coefficient characteristic.  They do not
bound the multiplicities of arbitrary nonabelian
monolithic crowns.  The fixed-rank representation
counts and the bounded-conductor product examples
elsewhere in this paper are kept separate from this
remaining full finite-generation assertion.
\end{remark}
